\documentclass[12pt,a4paper]{article}
\usepackage{graphicx}
\usepackage{amsmath}
\usepackage{hyperref}
\usepackage{geometry}
\geometry{margin=2.5cm}

\title{\textbf{CDFD Runtime Paper 08: Proto-Regime, Boundary Conditions, and Theoretical Discovery}}
\author{Steve Bico Mujjabi, MD\\
\small Independent Researcher\\
\small ORCID: \href{https://orcid.org/0009-0001-0556-5516}{0009-0001-0556-5516}}
\date{May 2026}

\begin{document}
\maketitle

\begin{abstract}
\noindent \textbf{Executive synthesis.} The Proto-Regime is best understood as
a boundary-condition discovery candidate: a high-flow, low-constraint state in
which ordinary interpretation of $\Psi_s$ breaks down. Earlier drafts described
it too strongly. This upgraded paper frames the Proto-Regime as a reproducible
runtime stress test and a theoretical hypothesis, not as proof of faster-than-
light travel, pre-Big-Bang physics, or deployable vacuum engineering.
\end{abstract}

\begin{figure}[ht]
\centering
\fbox{\begin{minipage}{0.92\textwidth}
\centering
\textbf{Runtime evidence path}\\[0.4em]
Prior CDFD mathematics $\longrightarrow$ CDFL/runtime state $\longrightarrow$ bounded output $\longrightarrow$ falsification target.
\end{minipage}}
\caption{Conceptual schematic for the runtime papers. The engine translates the mathematical frame from the prior CDFD papers into executable CDFL/runtime diagnostics; candidate outputs remain tied to explicit tests before any stronger claim is made.}
\end{figure}

\section{Prior CDFD Basis}
This manuscript does not rederive the mathematical core of CDFD. It uses the
Mujjabi Adaptive Operating Ratio and the constrained-vacuum notation introduced
in Part I \cite{MujjabiI2026}, the torus-knot hierarchy and boundary-mode work
from the Part I sequence \cite{MujjabiVI2026,MujjabiIX2026}, the public balance
law developed in the Part I capstone \cite{MujjabiX2026}, the Part I transport
tests \cite{MujjabiXII2026}, and the Life Number and tri-regime biology bridge
from Part II \cite{MujjabiPartII2026}. The runtime papers therefore act as an
execution layer: they keep the mathematics connected to commands, outputs,
falsification tests, and domain-specific limits.


\section{Definition}
The operating ratio is
\begin{equation}
    \Psi_s=\frac{\Phi}{C}S M_s.
\end{equation}
The Proto-Regime refers to simulations where $\Phi$ is forced high while
$C$ is forced toward a numerical floor. In this limit, ordinary flow-through-
constraint interpretation becomes unstable because the denominator approaches
zero.

\section{Why It Matters}
A good runtime must know where its own equations stop behaving like ordinary
domain models. The Proto-Regime is one such boundary. It marks a region where
the model triggers warnings, finite-value checks, and explicit caveats.

\section{Experiment Surface}
The local script \texttt{experiments/discover\_proto\_regime.py} is the named
entrypoint for probing this boundary. Related discovery machinery appears in:
\begin{itemize}
    \item \texttt{experiments/notebooks/discovery\_frontier\_sweep.py};
    \item \texttt{experiments/reports/NEW\_DISCOVERIES\_REPORT.md};
    \item \texttt{experiments/reports/FRONTIER\_DISCOVERIES.md}.
\end{itemize}
These reports are simulation logs. A high novelty score or new
schema name means the engine found an unusual modeled state, not that nature has
accepted the interpretation.

\section{Mathematical Boundary}
Let $C=\epsilon$ with $\epsilon>0$ small. Then
\begin{equation}
    \Psi_s \sim \frac{\Phi}{\epsilon}S M_s.
\end{equation}
For fixed $\Phi$, $S$, and $M_s$, $\Psi_s$ diverges as $\epsilon \to 0$. The
question is not whether the number becomes large; that is expected. The real
question is whether the state produces reproducible structure under bounded
regularization, or whether it is merely a division-by-small-number artifact.

\section{Part I Boundary Marker}
The active Part I boundary marker used here is the Paper XII superconductivity
reallocation sample with transport index
\[
    \mathcal{S}_{\mathrm{SC}}=4.187183192107868
\]
at the row where \(T/T_c=0.5287891440501045\). That row appears in the Part I
Paper XII superconductivity reallocation table and belongs to the Part I
transport-test program \cite{MujjabiXII2026}. The runtime Proto-Regime is
therefore written as a boundary family anchored to the Part-I-derived
4.187183192107868 transport marker.

\section{Interpretation Rules}
The Proto-Regime may be used as:
\begin{itemize}
    \item a numerical stress test;
    \item a boundary marker for model validity;
    \item a candidate theoretical state for future mathematical analysis;
    \item a warning system for non-physical parameter choices.
\end{itemize}
It does not establish:
\begin{itemize}
    \item frictionless travel;
    \item physical vacuum manipulation;
    \item a device claim;
    \item a clinical or engineering recommendation.
\end{itemize}

\section{Falsification and Controls}
Controls include:
\begin{enumerate}
    \item increasing the constraint floor $\epsilon$ and checking whether the
    pattern vanishes;
    \item adding finite $S$ and $M_s$ caps;
    \item varying grid resolution and time step;
    \item rerunning with different seeds and boundary conditions;
    \item comparing against a known numerical singularity test.
\end{enumerate}

\section{Connection to the CLI}
The CLI now needs a safe boundary-test command. Until then,
Proto-Regime experiments remain Python scripts and report entries. Any future
CDFL version needs an explicit warning when $C$ is below a declared
physical or domain-specific floor.

\section{Conclusion}
The Proto-Regime is important precisely because it is dangerous to overread.
It marks a boundary where the runtime can teach its authors which claims need
more mathematics, more controls, and stronger validation before they become
physics.
\begin{thebibliography}{9}
\bibitem{MujjabiI2026}
Steve Bico Mujjabi, MD. \emph{Vortex Stability in a Constrained Transport Vacuum and the Origin of the Fine-Structure Constant}. CDFD Part I Paper I, preprint, 2026.

\bibitem{MujjabiVI2026}
Steve Bico Mujjabi, MD. \emph{The Universal Torus Knot Hierarchy: $Z_n$ Symmetry, the Cinquefoil Family, and the General Structure of CDFT Particle Families}. CDFD Part I Paper VI, preprint, 2026.

\bibitem{MujjabiIX2026}
Steve Bico Mujjabi, MD. \emph{Even-$n$ Torus Knots in CDFT: The Exclusion of $T(2,2)$, the Extension to $T(2,2m)$, and the Anti-Phase Mass Scaling Law}. CDFD Part I Paper IX, preprint, 2026.

\bibitem{MujjabiX2026}
Steve Bico Mujjabi, MD. \emph{The Vacuum Equation of State: Deriving Torus Knot Self-Consistency from the Public CDFD Balance Law $\Psi = \Phi/C$}. CDFD Part I Paper X, preprint, 2026.

\bibitem{MujjabiXII2026}
Steve Bico Mujjabi, MD. \emph{Friction, Superconductivity, Plasma States, and Transport Mysteries: Observable Tests of Adaptive Constraint}. CDFD Part I Paper XII, preprint, 2026.

\bibitem{MujjabiPartII2026}
Steve Bico Mujjabi, MD. \emph{CDFD Part II: Origins of Life and Tri-Regime Bioenergetics}. Public release archive, 2026, \url{https://doi.org/10.5281/zenodo.20250821}.
\end{thebibliography}
\end{document}
